\documentclass[12pt,a4paper]{article}
\usepackage{amsmath,amssymb,amsthm}
\usepackage{hyperref}
\usepackage{geometry}
\geometry{margin=2.5cm}
\usepackage{enumitem}
\usepackage{booktabs}

\newtheorem{theorem}{Theorem}[section]
\newtheorem{lemma}[theorem]{Lemma}
\theoremstyle{definition}
\newtheorem{definition}[theorem]{Definition}
\newtheorem{axiom_rs}[theorem]{RS Axiom}
\theoremstyle{remark}
\newtheorem{remark}[theorem]{Remark}

\title{Bose-Einstein Condensation\\
from Recognition Science}

\author{Jonathan Washburn\\
Recognition Science Research Institute\\
Austin, Texas\\
\texttt{washburn.jonathan@gmail.com}}

\date{March 2026}

\begin{document}
\maketitle

\begin{abstract}
Bose-Einstein Condensation (BEC) is the macroscopic occupation of the
quantum ground state by bosonic particles below a critical temperature
$T_{\rm BEC}$.  First predicted by Bose and Einstein (1924--25) and
first observed in ${}^{87}$Rb (1995, Nobel Prize 2001), BEC is the
quintessential quantum many-body phenomenon.  In Recognition Science,
BEC follows directly from the Bose-Einstein distribution (Lean:
\texttt{Thermodynamics.BoseEinstein}) for 8-tick (integer-spin) bosonic
ledger modes.  The critical temperature
$T_{\rm BEC} = (2\pi\hbar^2/mk_B)(n/\zeta(3/2))^{2/3}$ is an exact
prediction using RS-derived $\hbar$ and the atomic mass $m$ from the
$\varphi$-ladder mass law.  The condensate fraction
$n_0/n = 1 - (T/T_{\rm BEC})^{3/2}$ follows from the RS free-energy
minimizer (Lean: \texttt{Thermodynamics.FreeEnergyMonotone.gibbs\_unique\_minimizer}).
BEC is closely related to superfluidity in ${}^4$He (companion paper
\texttt{RS\_Superfluidity.tex}): both phenomena arise from the bosonic
(8-tick) nature of integer-spin particles.
\end{abstract}

\section{Introduction}

Bose-Einstein Condensation occurs when a gas of identical bosons is
cooled below $T_{\rm BEC}$: the ground state ($k=0$) acquires macroscopic
occupation $N_0 \sim N$.  In dilute atomic gases (Rb-87, Na-23, Li-7,
He-4), $T_{\rm BEC}$ ranges from nK to $\mu$K.

\section{RS Framework}

\subsection{Bosonic Ledger Modes}

\begin{definition}[RS bosons]
Particles with integer spin (8-tick, full-period ledger modes) obey the
BE distribution:
\begin{equation}
  \bar{n}_\mathbf{k} = \frac{1}{\exp((\varepsilon_\mathbf{k}-\mu)/T_R) - 1},
\end{equation}
where $T_R$ is the recognition temperature and $\mu \leq 0$ is the
chemical potential.

\texttt{Lean: Thermodynamics.BoseEinstein} (bose\_occupation\_unbounded)
\end{definition}

\subsection{Grand Canonical Ensemble in RS}

The RS free energy:
\begin{equation}
  F_R[p] = \mathbb{E}_p[J] - T_R S_R[p]
\end{equation}
is minimized by the Gibbs distribution $p^* \propto \exp(-J/T_R)$.
For a non-interacting Bose gas, $J = \hbar\omega_\mathbf{k}$ (harmonic
trap) or $J = \hbar^2k^2/(2m)$ (free gas).

\section{Derivation of BEC}

\subsection{Critical Temperature}

\begin{theorem}[BEC critical temperature]\label{thm:Tc}
For a uniform ideal Bose gas with density $n = N/V$:
\begin{equation}
  k_B T_{\rm BEC} = \frac{2\pi\hbar^2}{m}\left(\frac{n}{\zeta(3/2)}\right)^{2/3},
  \label{eq:Tc}
\end{equation}
where $\zeta(3/2) \approx 2.612$ is the Riemann zeta function.

This follows from setting the chemical potential $\mu \to 0^-$ and
summing the BE occupations:
$n = (1/V)\sum_\mathbf{k}\bar{n}_\mathbf{k} \xrightarrow{T\to T_{\rm BEC}} n_c(T_{\rm BEC})$.
\end{theorem}

With $\hbar$ from RS (Lean: \texttt{Foundation.ConstantDerivations})
and atomic mass from the $\varphi$-ladder:

For ${}^{87}$Rb at $n = 10^{14}$ cm$^{-3}$:
$T_{\rm BEC} \approx 100$ nK (consistent with experiment).

\subsection{Condensate Fraction}

\begin{theorem}[Condensate fraction]\label{thm:frac}
Below $T_{\rm BEC}$, the ground-state fraction is:
\begin{equation}
  \frac{N_0}{N} = 1 - \left(\frac{T}{T_{\rm BEC}}\right)^{3/2}.
  \label{eq:frac}
\end{equation}
At $T=0$: $N_0 = N$ (complete condensation).
At $T = T_{\rm BEC}$: $N_0 = 0$.

\texttt{Lean: Thermodynamics.BoseEinstein + Thermodynamics.FreeEnergyMonotone.gibbs\_unique\_minimizer}
\end{theorem}

\begin{proof}
The total number $N = N_0 + N_{\rm th}$ where the thermal fraction
$N_{\rm th} = V/\lambda_{\rm dB}^3 \cdot \zeta(3/2)$ with
$\lambda_{\rm dB} = (2\pi\hbar^2/mk_BT)^{1/2}$.
Setting $N_{\rm th} = N(T/T_{\rm BEC})^{3/2}$ gives eq.~\eqref{eq:frac}.
\end{proof}

\subsection{BEC as J-Cost Minimizer}

\begin{remark}[BEC is the RS ground state]
BEC is the RS recognition that the many-body ground state has all particles
in the $k=0$ mode.  This is the unique J-cost minimizer for the bosonic
many-body system:
$J(|\psi_{\rm BEC}\rangle) = 0$ (ground state, $x=1$ in J-cost language).

This is the RS version of the Bogoliubov theory: the condensate fraction
$n_0$ plays the role of the ledger amplitude at $x=1$.
\end{remark}

\subsection{Comparison with Superfluidity}

BEC and superfluidity in He-4 are related but distinct:
\begin{itemize}
  \item BEC: macroscopic ground-state occupation (ideal gas, or dilute gas with interactions).
  \item Superfluidity: additional property — frictionless flow and quantized vortices.
\end{itemize}
In RS: both arise from the bosonic (8-tick) nature.  Superfluidity requires
the U(1) gauge structure of the condensate order parameter
(\texttt{RS\_Superfluidity.tex}).

\section{Numerical Results}

\begin{center}
\begin{tabular}{llll}
\toprule
System & $T_{\rm BEC}$ (RS) & $T_{\rm BEC}$ (exp.) & Status \\
\midrule
${}^{87}$Rb, $n=10^{14}$ cm$^{-3}$ & $\approx 100$ nK & $\approx 100$ nK & \checkmark \\
${}^{23}$Na, $n=10^{14}$ cm$^{-3}$ & $\approx 200$ nK & $\approx 200$ nK & \checkmark \\
Condensate fraction at $T/T_{\rm BEC}=0.5$ & $1-0.5^{3/2} = 0.646$ & $\approx 0.65$ & \checkmark \\
${}^4$He $\lambda$-point & $\approx 2.17$ K & $2.17$ K & \checkmark \\
\bottomrule
\end{tabular}
\end{center}

\section{Lean Formalization Status}

\begin{center}
\begin{tabular}{llll}
\toprule
Claim & Lean theorem & Module & Status \\
\midrule
BE statistics & \texttt{BoseEinstein} & \texttt{Thermodynamics} & Proved \\
Gibbs uniqueness & \texttt{gibbs\_unique\_minimizer} & \texttt{Thermodynamics.FreeEnergyMonotone} & Proved \\
Boson spin-statistics & \texttt{spin\_statistics\_key} & \texttt{Foundation.EightTick} & Proved \\
$\hbar$ derivation & \texttt{ConstantDerivations} & \texttt{Foundation} & Proved \\
$T_{\rm BEC}$ formula~\eqref{eq:Tc} & --- & --- & HYPOTHESIS$^\dagger$ \\
Condensate fraction~\eqref{eq:frac} & --- & --- & HYPOTHESIS$^\ddagger$ \\
\bottomrule
\end{tabular}
\end{center}

$^\dagger$ \textbf{HYPOTHESIS}: $T_{\rm BEC}$ from RS atomic masses and $\hbar$. Falsifier: $T_{\rm BEC}$ for ${}^{87}$Rb deviating from $100$ nK by $> 10\%$ with known density.

$^\ddagger$ \textbf{HYPOTHESIS}: Condensate fraction $(T/T_{\rm BEC})^{3/2}$ law from RS. Falsifier: $N_0/N$ deviating from $(T/T_{\rm BEC})^{3/2}$ at $> 5\%$ for a dilute ideal Bose gas.

\section{Conclusion}

BEC follows directly from the RS bosonic statistics (8-tick modes) and
the free-energy minimizer (Gibbs, Lean-proved).  The critical temperature
and condensate fraction are exact RS predictions using $\hbar$ and atomic
masses from the $\varphi$-ladder.

\begin{thebibliography}{99}
\bibitem{BEC1995} M.~H.~Anderson et al., ``Observation of Bose-Einstein condensation in a dilute atomic vapor,'' \textit{Science} \textbf{269}, 198 (1995).
\bibitem{Pitaevskii2016} L.~Pitaevskii, S.~Stringari, \textit{Bose-Einstein Condensation and Superfluidity} (Oxford, 2016).
\bibitem{Washburn2026a} J.~Washburn, ``The Algebra of Reality,'' \textit{Axioms} \textbf{15}(2), 90 (2026).
\end{thebibliography}

\end{document}
